% =============================================================================
%  lec01.tex —— 第 1 讲 引言（底本 PDF 第 1–7 页）
% =============================================================================
\lecture{1}{引言}
\begin{center}\small\color{black!55}（底本 PDF 第 1–7 页）\end{center}

本课程的目标，是对 $\infty$-范畴理论——或者更一般地，对同伦连贯的数学——作一个不重技术的介绍。课程侧重想法与动机，希望能充当这一领域基础参考文献的向导，包括 Lurie 的几本书 [Lur09]、[Lur12]、[Lur]。

\section{一个口号}

\begin{slogan}\label{L1-1.1}\centering\par
$\infty$-范畴论 $=$ 范畴论 $+$ 同伦论。
\end{slogan}

\section{经典范畴论}

\nitem\label{L1-2.1} 范畴由 Eilenberg 与 MacLane 于 1945 年在他们关于代数拓扑与同调代数的研究中引入。

\begin{definition}\label{L1-2.2}
一个\textbf{范畴} $\cC$ 由以下数据组成：
\begin{itemize}[leftmargin=2.2em,itemsep=0.1em]
\item 一个类 $\Ob(\cC)$，其元素称为\textbf{对象}；
\item 对任意两个对象 $a$ 与 $b$，一个类 $\Hom(a,b)$，其元素称为从 $a$ 到 $b$ 的\textbf{态射}，记作 $f\colon a\to b$；
\item 一个二元运算 $\circ$，称为态射的\textbf{复合}，使得对任意三个对象 $a$、$b$、$c$，有映射
\[
\circ\colon \Hom(a,b)\times\Hom(b,c)\to \Hom(a,c).
\]
\end{itemize}
上述数据应满足两条公理：
\begin{itemize}[leftmargin=2.2em,itemsep=0.1em]
\item 结合律：对态射 $f\colon a\to b$、$g\colon b\to c$ 与 $h\colon c\to d$，
\[
h\circ(g\circ f)=(h\circ g)\circ f.
\]
\item 单位律：对任意对象 $x$，存在态射 $\id_x\colon x\to x$，称为 $x$ 的\textbf{单位态射}，使得对任意态射 $f\colon a\to b$，有
\[
\id_b\circ f=f=f\circ \id_a.
\]
\end{itemize}
\end{definition}

\nitem\label{L1-2.3} 范畴论的力量体现在下述原则中：
\begin{center}
\emph{要研究一族对象，就还应当考虑这些对象之间适当定义的态射。}
\end{center}
人们把自己喜爱的数学实体——往往是带结构的集合——组装成一个范畴，并把态射宣告为保持这些结构的函数。例子包括 $\Set$、$\Grp$、$\Ring$、$\Top$ 等。

\nitem\label{L1-2.4} 我们想强调（经典）范畴论中的下述几条信条：
\begin{enumerate}[label=(\arabic*),leftmargin=2.4em,itemsep=0.1em]
\item 态射是离散的：对两个态射 $f$ 与 $g$，人们可以说 $f=g$ 或 $f\neq g$，而这是所能作出的唯一比较。
\item 结合律是严格的：$h\circ(g\circ f)$ 与 $(h\circ g)\circ f$ 被要求相等，而不是在较弱意义下等价。
\item 复合是具体的：$g\circ f$ 没有歧义。
\end{enumerate}
本课程将把这三条全部放弃。

\section{走向高阶范畴}

\nitem\label{L1-3.1} 我们将放弃信条 (1)，并给 $\Hom(a,b)$ 赋予更丰富的结构：它们可以是拓扑空间，甚至可以是范畴本身。后者定义了\textbf{严格 $2$-范畴}：它有对象与态射，还有态射之间的态射，称为 \textbf{$2$-态射}。$2$-态射的复合有两种：\textbf{纵复合}与\textbf{横复合}。

\yizhu{原书此处有一幅示意图，画出 $2$-态射的纵复合与横复合。图形本身不含文字信息，译本从略，见底本第 2 页。}

这些复合应满足一列结合律与单位律公理，而它们全部用\textbf{等号}陈述。

归纳下去，就得到严格 $n$-范畴的概念。用经典范畴论的语言，我们有：

\begin{definition}\label{L1-3.2}
\textbf{严格 $n$-范畴}是以严格 $(n-1)$-范畴充实的范畴，即其 $\Hom$ 对象本身带有严格 $(n-1)$-范畴结构的范畴。
\end{definition}

\nitem\label{L1-3.3} 然而，严格 $n$-范畴的有趣例子极少：

\begin{example}\label{L1-3.4}
只有一个对象 $\ast$ 的严格 $2$-范畴，等同于如下数据：一个范畴 $\cC\coloneqq \Hom(\ast,\ast)$，配以一个乘法函子 $-\otimes -\colon \cC\times\cC\to\cC$，它严格结合且严格单位。这个定义是\textbf{邪恶的}\footnote{同构原则：关于固定范畴中对象的一切语法正确的性质，都应在同构下不变。}且\textbf{不切实际的}\footnote{即便对集合，$(X\times Y)\times Z=X\times(Y\times Z)$ 在 Zermelo--Fraenkel 公理系统中也没有意义。}。正确的概念是\textbf{幺半范畴}：我们改为提供自然同构
\[
X\otimes(Y\otimes Z)\xrightarrow{\ \simeq\ }(X\otimes Y)\otimes Z,\qquad
\mathbf{1}\otimes X\xrightarrow{\ \simeq\ }X\xleftarrow{\ \simeq\ }X\otimes \mathbf{1},
\]
并要求它们满足某些连贯性条件。
\end{example}

\nitem\label{L1-3.5} 上例提示我们还应放弃信条 (2)，允许结合律在较弱的意义下成立。这引向弱 $n$-范畴的概念——注意只是概念，还不是定义——当 $n=\infty$ 时则引向弱 $\omega$-范畴。

对于弱 $n$-范畴的定义，我们至少有如下愿望：
\begin{itemize}[leftmargin=2.2em,itemsep=0.1em]
\item 弱 $1$-范畴就是范畴。
\item 只有一个对象的弱 $2$-范畴，等同于一个幺半范畴的数据。
\item 对弱 $n$-范畴中任意两个对象 $a$ 与 $b$，$\Hom(a,b)$ 应当是弱 $(n-1)$-范畴。
\item 其定义应满足同构原则。
\end{itemize}

把最后两条愿望合起来，我们得到：
\begin{center}
\emph{当 $n\ge 2$ 时，我们绝不应当要求弱 $n$-范畴中的两个态射相等。}
\end{center}
再由归纳，
\begin{center}
\emph{当 $n>k$ 时，我们绝不应当要求弱 $n$-范畴中的两个 $k$-态射相等。}
\end{center}
这些看似无害的愿望会导致一场组合爆炸的噩梦。

\nitem\label{L1-3.6} 设 $f$、$g$、$h$ 是弱 $n$-范畴中可复合的态射。按前面的讨论，在结合律公理中，$h\circ(g\circ f)$ 与 $(h\circ g)\circ f$ 应当等价而不是相等。

然而，正如幺半范畴的定义所提示的，这个等价必须被看作\textbf{结构}而不是性质：我们需要提供从 $h\circ(g\circ f)$ 到 $(h\circ g)\circ f$ 的一个可逆 $2$-态射，称为\textbf{结合子}，使得下图交换：
\begin{center}
\begin{tikzcd}[column sep=3.2em,row sep=2.6em]
(kh)(gf) \arrow[rr] \arrow[dd] & & k(h(gf)) \arrow[dd] \\
& & \\
((kh)g)f \arrow[rr] & & k((hg)f)
\end{tikzcd}
\end{center}
然而，当 $n\ge 3$ 时，连 $2$-态射的这种交换性也应理解为等价而非相等，并应由一个可逆 $3$-态射来见证，比如说，从顺时针弧到逆时针弧的那个 $3$-态射。这些 $3$-态射本身又应使某个图交换，而当 $n\ge 4$ 时，这个交换性又由可逆 $4$-态射见证……

更糟的是，我们还必须处理高阶态射复合的结合律公理，而这类公理有多种类型。提醒一句：$2$-态射既能纵复合也能横复合，后者还与 $1$-态射的复合相互作用。

这一无穷无尽的结合律清单，称为\textbf{连贯数据}，很快就变得无法写出，也难以操作。

\nitem\label{L1-3.7} 这场痛苦地以组合方式定义弱 $n$-范畴的追求，由 B\'enabou 于 1967 年开启，大约在 2000 年代初终止。

\nitem\label{L1-3.8} 拯救高阶范畴学家（以及本课程）的，是 Grothendieck 在 1983 年前后所写的《探究叠》（\emph{Pursuing Stacks}）中提出的下述同伦假设：

\begin{slogan}\label{L1-3.9}\centering\par
$\infty$-群胚的理论应当与空间的同伦论相同。
\end{slogan}

\nitem\label{L1-3.10} 这里 $\infty$-群胚指所有态射与高阶态射都可逆的弱 $\omega$-范畴。一般地，$(n,k)$-范畴指 $m>k$ 时的 $m$-态射都可逆的弱 $n$-范畴。当 $n=\infty$ 时，上述弱 $\infty$-范畴就是弱 $\omega$-范畴\footnote{$(\infty,1)$-范畴通常就简称为 $\infty$-范畴。}。

洞见在于：结合律的连贯数据在组合上令人生畏，却可以隐藏在空间的同伦论中。

这给出发展高阶范畴论的一条路径：我们首先宣告 $\infty$-群胚——即 $(\infty,0)$-范畴——就是空间的同伦型，然后归纳地定义 $(\infty,k)$-范畴。在归纳步中，关于可逆高阶态射的结合律连贯数据，已经由 $\infty$-群胚的理论驯服。

\nitem\label{L1-3.11} 人们可能会问：既然 $(\infty,0)$-范畴的理论等同于拓扑空间的同伦论，那么 $(\infty,1)$-范畴的理论是否应当等同于拓扑范畴——即以拓扑空间充实的范畴——的同伦论？

答案是：是的，但我们必须先理解后者的含义。这需要 Quillen 的抽象同伦论，即模型范畴，它是下一讲的内容。

不过，拓扑范畴并不是 $(\infty,1)$-范畴最方便的模型，至少在后者的某些有趣应用中如此。目前发展最完善的模型，得益于 Lurie 的几本书，是\textbf{拟范畴}；正如我们暗示过的，它放弃信条 (3)。

\section{同伦假设}

\nitem\label{L1-4.1} 同伦论可追溯到 Poincar\'e 自 1895 年起关于基本群的工作。

\begin{definition}\label{L1-4.2}
设 $X$、$Y$ 是拓扑空间，$f,g\colon X\to Y$ 是连续函数。$f$ 与 $g$ 之间的一个\textbf{同伦}是连续函数 $H\colon X\times[0,1]\to Y$，使得 $H(-,0)=f$ 且 $H(-,1)=g$。若 $f$ 与 $g$ 之间存在同伦，则称它们\textbf{同伦}。
\end{definition}

设 $X$、$Y$ 是拓扑空间。$X$ 与 $Y$ 之间的一个\textbf{同伦等价}是一对连续函数 $p\colon X\to Y$ 与 $q\colon Y\to X$，使得 $q\circ p$ 同伦于 $\id_X$，且 $p\circ q$ 同伦于 $\id_Y$。

\nitem\label{L1-4.3} 经典同伦论关注拓扑空间在同伦等价下不变的信息，例如其同伦群。这样的信息可以编码进一个范畴。

\begin{construction}\label{L1-4.4}
\textbf{同伦范畴} $\hTop$ 定义如下：
\begin{itemize}[leftmargin=2.2em,itemsep=0.1em]
\item 对象是性质良好的拓扑空间，例如 CW 复形；
\item 态射是连续映射的同伦类，态射的复合由连续映射的复合诱导。
\end{itemize}
\end{construction}

\begin{variant}\label{L1-4.5}
类似地，考虑带基点的空间，即配有基点的空间，我们得到范畴 $\hTop_\ast$。
\end{variant}

\nitem\label{L1-4.6} 按设计，同伦等价恰好就是 $\hTop$ 中的等价。事实上，$\hTop$ 可由性质良好的拓扑空间范畴将同伦等价形式地可逆化得到；它也可由全体拓扑空间的范畴 $\Top$ 将弱同伦等价\footnote{连续映射是\textbf{弱同伦等价}，是指它诱导各 $\pi_k$ 之间的同构。由 Whitehead 的一个定理，CW 复形之间的弱同伦等价是同伦等价。}形式地可逆化得到。

\begin{definition}\label{L1-4.7}
\textbf{同伦型}是 $\hTop$ 中的一个对象 $X$。若对任意基点，当 $k>n$ 时 $\pi_k(X,x)\cong 0$，则称 $X$ 是\textbf{同伦 $n$-型}。
\end{definition}

设 $\hTop_{\le n}\subset\hTop$ 是由同伦 $n$-型构成的全子范畴。

\nitem\label{L1-4.8} 事实证明，$\pi_0(X)$ 与 $\pi_1(X)$ 的信息可以由一个与 $X$ 相伴的范畴捕捉。

\begin{construction}\label{L1-4.9}
设 $X$ 是拓扑空间。$X$ 的\textbf{基本群胚} $\pi_{\le 1}X$ 是如下定义的范畴：
\begin{itemize}[leftmargin=2.2em,itemsep=0.1em]
\item 对象是 $X$ 的点；
\item 态射是 $X$ 中道路的同伦类，复合由区间的拼接诱导。
\end{itemize}
\end{construction}

\begin{exercise}\label{L1-4.10}
验证上述构造满足结合律公理。
\end{exercise}

\nitem\label{L1-4.11} 容易看出：
\begin{itemize}[leftmargin=2.2em,itemsep=0.1em]
\item 基本群胚 $\pi_{\le 1}X$ 确实是群胚，即所有态射都可逆；
\item 集合 $\pi_0(X)$ 可与 $\pi_{\le 1}X$ 中对象的同构类集合等同；
\item 有自然的群同态 $\pi_1(X,x)\cong\Hom(x,x)$，其中右端的 $x$ 视为 $\pi_{\le 1}X$ 中的对象。
\end{itemize}

\begin{exercise}\label{L1-4.12}
证明 $\pi_{\le 1}$ 定义了从 $\hTop_{\le 1}$ 到小群胚范畴 $\hGrpd$ 的等价，其中态射由函子的等价类给出。提示：Eilenberg--MacLane 空间。
\end{exercise}

\nitem\label{L1-4.13} 受上述结果鼓舞，人们可能尝试如下构造一个 $2$-范畴 $\pi_{\le 2}X$：
\begin{itemize}[leftmargin=2.2em,itemsep=0.1em]
\item 对象是 $X$ 的点；
\item 态射是 $X$ 中的道路，复合由区间的拼接诱导；
\item $2$-态射是 $X$ 中道路之间同伦的同伦类，复合由方形的拼接诱导。
\end{itemize}
如果这个定义可行，注意：
\begin{itemize}[leftmargin=2.2em,itemsep=0.1em]
\item 所有态射与 $2$-态射都可逆；
\item 集合 $\pi_0(X)$ 可与 $\pi_{\le 2}X$ 中对象的同构类集合等同；
\item 群 $\pi_1(X,x)$ 可与 $\Hom(x,x)$ 中对象的同构类集合等同；
\item 有自然的群同态 $\pi_2(X,x)\cong\Hom(\id_x,\id_x)$，其中右端的 $\id_x$ 视为 $\Hom(x,x)$ 中的对象。
\end{itemize}

\begin{exercise}\label{L1-4.14}
用 $\pi_{\le 2}X$ 描述 $\pi_1(X,x)$ 在 $\pi_2(X,x)$ 上的作用。
\end{exercise}

\begin{exercise}\label{L1-4.15}
在上述 $\pi_{\le 2}X$ 的定义中，我们还能把态射定义为 $X$ 中道路的同伦类吗？请说服自己：那样 $2$-态射将没有良好定义。

如果卡住了，试试下面这一点：在 $\pi_{\le 1}X$ 的定义中，我们能把对象定义为点的同伦类，即 $X$ 的连通分支吗？
\end{exercise}

\nitem\label{L1-4.16} 注意 $\pi_{\le 2}X$ 不可能是严格的：对可复合的道路 $f,g,h\colon[0,1]\to X$，复合
\begin{equation}\label{L1-4.1eq}
(h\circ(g\circ f))(t)=
\begin{cases}
f(4t), & t\in[0,1/4],\\
g(4t-1), & t\in[1/4,1/2],\\
h(2t-1), & t\in[1/2,1],
\end{cases}
\end{equation}
与
\begin{equation}\label{L1-4.2eq}
((h\circ g)\circ f)(t)=
\begin{cases}
f(2t), & t\in[0,1/2],\\
g(4t-2), & t\in[1/2,3/4],\\
h(4t-3), & t\in[3/4,1],
\end{cases}
\end{equation}
同伦但不等。此外，很难找到它们之间自然的同伦，尽管只要构造对任意 $X$ 都有效，所有这样的同伦彼此同伦。

最后这句话提供了 $\pi_{\le 2}X$ 中结合律的连贯数据。可以验证 $\pi_{\le 2}X$ 确实是弱 $2$-范畴。

\begin{challenge}\label{L1-4.17}
证明 $\pi_{\le 2}$ 定义了从 $\hTop_{\le 2}$ 到小弱 $2$-群胚范畴 $\hGrpd_2$ 的等价，其中态射由函子的等价类给出。
\end{challenge}

\nitem\label{L1-4.18} 当 $n\ge 3$ 时，要构造 $\pi_{\le n}X$，按练习 4.15，$2$-态射必须是道路之间的同伦，而不能是这种同伦的同伦类。因此，只要弱 $n$-范畴的理论以组合方式发展，我们就必须为从 \eqref{L1-4.1eq} 到 \eqref{L1-4.2eq} 选择一个同伦。这样的选择不自然，而我们还必须把它们塞进弱 $n$-范畴定义中结合律的连贯数据里。随着 $\pi_{\le n}X$ 中的 $n$ 增大，这样的同伦连贯数据变得令人生畏。

\nitem\label{L1-4.19} 启发式地说，我们有两项不可能的任务：
\begin{itemize}[leftmargin=2.2em,itemsep=0.1em]
\item 给出（弱）$n$-群胚的组合定义；
\item 验证——或者更确切地说，提供同伦连贯数据以使得——$\pi_{\le n}X$ 成为弱 $n$-群胚。
\end{itemize}
Grothendieck 的同伦假设说，这两项任务实际上是同一项任务，而我们两项都不该做。

\begin{slogan}\label{L1-4.20}\centering\par
$n$-群胚 $=$ 同伦 $n$-型；\qquad $\infty$-群胚 $=$ 同伦型。
\end{slogan}

\nitem\label{L1-4.21} 范畴论与同伦论的这次重聚甚至可以追溯到 Kan 在 1950 年代的工作，它是本课程的指导哲学。

\appendixsection{论严格与弱}

\begin{exercise}\label{L1-A.1}
证明任意弱 $2$-范畴都等价于一个严格 $2$-范畴。
\end{exercise}

\begin{exercise}\label{L1-A.2}
证明：
\begin{enumerate}[label=(\arabic*),leftmargin=2.4em,itemsep=0.1em]
\item 知道一个只有一个对象、一个态射的弱 $3$-范畴 $\cC$，等同于知道一个辫幺半范畴 $\cD$；
\item 在上述情形中，$\cC$ 等价于严格 $3$-范畴，当且仅当 $\cD$ 是对称幺半范畴。
\end{enumerate}
\end{exercise}

\begin{exercise}\label{L1-A.3}
证明 $\pi_{\le 3}S^2$ 不同价于严格 $3$-群胚。
\end{exercise}

\nitem\label{L1-A.4} \textbf{建议阅读}\quad [Sim98]。

\subsection*{参考文献}
\par\noindent\hangindent=5.2em\hangafter=1 [Lur]\hspace{0.6em}Jacob Lurie. \emph{Kerodon}. \url{https://kerodon.net}.
\par\noindent\hangindent=5.2em\hangafter=1 [Lur09]\hspace{0.6em}Jacob Lurie. \emph{Higher Topos Theory}. Princeton University Press, 2009.
\par\noindent\hangindent=5.2em\hangafter=1 [Lur12]\hspace{0.6em}Jacob Lurie. \emph{Higher Algebra}, 2012.
\par\noindent\hangindent=5.2em\hangafter=1 [Sim98]\hspace{0.6em}Carlos Simpson. Homotopy types of strict 3-groupoids. arXiv preprint \texttt{math/9810059}, 1998.
